\chapter{类、封装和不变量}

\section{问题入口：温度范围}

类不是把几个变量包起来。类的价值在于维护不变量。\term{不变量}{invariant} 指对象在任何对外可见状态下都必须成立的条件。假设要表示一个温度区间，要求最低温不能高于最高温。

朴素结构体写法如下：

\begin{lstlisting}[language=C++,caption={没有保护不变量的结构体}]
struct TemperatureRange {
    double low = 0.0;
    double high = 0.0;
};

TemperatureRange range{.low = 30.0, .high = 10.0}; // 逻辑错误
\end{lstlisting}

编译器不知道业务规则，所以这段代码能编译。要把规则放进类型里。

\section{构造函数建立有效状态}

\begin{lstlisting}[language=C++,caption={构造函数检查不变量}]
#include <stdexcept>

class TemperatureRange {
public:
    TemperatureRange(double low, double high)
        : low_{low}, high_{high}
    {
        if (high < low) {
            throw std::invalid_argument{"high must not be lower than low"};
        }
    }

    [[nodiscard]] double low() const noexcept
    {
        return this->low_;
    }

    [[nodiscard]] double high() const noexcept
    {
        return this->high_;
    }

    [[nodiscard]] bool contains(double value) const noexcept
    {
        return this->low_ <= value && value <= this->high_;
    }

private:
    double low_ = 0.0;
    double high_ = 0.0;
};
\end{lstlisting}

数据成员是私有的，调用者不能绕过构造函数随意改。成员函数里使用 \code{this->}，目的是把“当前对象的成员”写清楚。

\section{setter 不一定是封装}

很多人把封装理解成“成员私有，再给每个成员写 getter 和 setter”。这只是形式。如果写出下面的接口，不变量又被破坏了：

\begin{lstlisting}[language=C++,caption={危险的 setter}]
void set_low(double value)
{
    this->low_ = value; // 可能让 low_ > high_
}
\end{lstlisting}

更好的接口应该表达业务动作。例如扩展范围：

\begin{lstlisting}[language=C++,caption={用业务动作维护不变量}]
void expand_to_include(double value) noexcept
{
    if (value < this->low_) {
        this->low_ = value;
    }
    if (value > this->high_) {
        this->high_ = value;
    }
}
\end{lstlisting}

调用者不需要知道内部如何保存，只知道对象会把某个值纳入范围。

\section{值语义优先}

\code{TemperatureRange} 可以复制、赋值、放进 \code{std::vector}。这是值语义。值语义让对象像数字和字符串一样好用。

\begin{lstlisting}[language=C++,caption={值语义让对象易组合}]
std::vector<TemperatureRange> ranges;
ranges.emplace_back(10.0, 20.0);
ranges.emplace_back(-5.0, 8.0);
\end{lstlisting}

不要过早把所有类都设计成继承体系。许多业务概念更适合小而明确的值对象。

\section{本章边界检查}

设计类时至少问：

\begin{itemize}
  \item 对象什么时候有效？
  \item 哪些状态绝对不允许出现？
  \item 构造函数能否一次建立有效状态？
  \item 对外接口是业务动作，还是裸字段修改？
  \item 对象是否应该支持复制和移动？
\end{itemize}

\begin{keyidea}
类的核心不是语法里的 \code{class}，而是把数据和规则放在同一个边界内。好的类让错误状态难以构造，让正确用法自然出现。
\end{keyidea}

\section{从坏对象推导构造边界}

继续看温度区间。为什么不能让调用者先默认构造，再分别设置上下界？因为中间状态可能已经被观察到：

\begin{lstlisting}[language=C++,caption={分步设置制造临时坏状态}]
TemperatureRange range;
range.set_high(10.0);
range.set_low(30.0); // 如果允许，就破坏 low <= high
if (range.contains(20.0)) {
    // 这里的判断已经没有可信含义
}
\end{lstlisting}

这个反例推出一条类设计规则：如果一个对象必须同时拥有几项数据才有效，就让构造函数一次收齐这些数据，并在构造边界检查。不应该把“不完整但暂时别用”的对象暴露给普通调用者。

\section{成员函数应该维护对象语义}

假设业务需要把两个温度区间合并。不要把内部字段拿出来让调用者自己算，而是提供表达业务的函数：

\begin{lstlisting}[language=C++,caption={合并两个有效区间}]
TemperatureRange merge(TemperatureRange lhs, TemperatureRange rhs)
{
    const double low = std::min(lhs.low(), rhs.low());
    const double high = std::max(lhs.high(), rhs.high());
    return TemperatureRange{low, high};
}
\end{lstlisting}

因为 \code{lhs} 和 \code{rhs} 都已经是有效对象，\code{min(low)} 和 \code{max(high)} 构造出来的新对象也应当有效。这里的推理链条很重要：类不变量让函数可以建立在更少的防御性检查上。

\section{什么时候用 struct，什么时候用 class}

\cpp 里 \code{struct} 和 \code{class} 的主要语法差异是默认访问权限；真正的设计差异来自意图。只承载数据、没有复杂不变量的聚合，适合 \code{struct}：

\begin{lstlisting}[language=C++,caption={简单结果适合 struct}]
struct ScoreSummary {
    std::size_t count = 0;
    int sum = 0;
    double average = 0.0;
};
\end{lstlisting}

需要保护规则、隐藏表示、控制修改路径的概念，适合 \code{class}。这不是硬性宗教，而是让读者一看到类型形态，就能大致知道它的使用方式。

\section{类的测试不是只测 getter}

测试 \code{TemperatureRange} 时，不要只检查 \code{low()} 和 \code{high()} 返回原值。更重要的是不变量：

\begin{lstlisting}[language=C++,caption={围绕不变量测试类}]
const TemperatureRange normal{10.0, 20.0};
assert(normal.contains(10.0));
assert(normal.contains(15.0));
assert(!normal.contains(21.0));

bool rejected = false;
try {
    const TemperatureRange invalid{30.0, 10.0};
} catch (const std::invalid_argument&) {
    rejected = true;
}
assert(rejected);
\end{lstlisting}

如果一个类的测试只是在验证字段保存和取出，说明这个类可能没有承担真正职责；它也许只是一个 \code{struct}，或者需要重新思考业务动作。

\begin{exercisebox}
实现一个 \code{Date} 类，要求月份在 \code{1} 到 \code{12}，日期必须符合月份天数。先不处理闰年，再把闰年规则加进去。练习重点不是背日期算法，而是把“不允许出现的状态”挡在构造函数外。
\end{exercisebox}

\section{Date 案例：从非法状态推导类}

温度区间只有两个字段，规则也简单。日期更接近真实业务：年、月、日三个字段一起决定对象是否有效。朴素结构体很容易制造坏对象：

\begin{lstlisting}[language=C++,caption={日期坏对象}]
struct DateData {
    int year = 0;
    int month = 0;
    int day = 0;
};

DateData impossible{.year = 2026, .month = 2, .day = 31};
\end{lstlisting}

编译器不知道二月没有三十一日。于是日期不能长期只靠裸 \code{struct} 表达。先写规则函数：

\begin{lstlisting}[language=C++,caption={日期规则函数}]
constexpr bool is_leap_year(int year) noexcept
{
    return (year % 400 == 0) ||
           (year % 4 == 0 && year % 100 != 0);
}

constexpr int days_in_month(int year, int month) noexcept
{
    constexpr int DAYS_IN_MONTH[] = {
        0, 31, 28, 31, 30, 31, 30,
        31, 31, 30, 31, 30, 31,
    };
    if (month == 2 && is_leap_year(year)) {
        return 29;
    }
    return DAYS_IN_MONTH[month];
}
\end{lstlisting}

规则函数单独存在，类负责调用它建立有效状态：

\begin{lstlisting}[language=C++,caption={Date 构造边界}]
class Date {
public:
    Date(int year, int month, int day)
        : year_{year}, month_{month}, day_{day}
    {
        if (month < 1 || month > 12) {
            throw std::invalid_argument{"month must be 1..12"};
        }
        if (day < 1 || day > days_in_month(year, month)) {
            throw std::invalid_argument{"day is invalid for month"};
        }
    }

    [[nodiscard]] int year() const noexcept { return this->year_; }
    [[nodiscard]] int month() const noexcept { return this->month_; }
    [[nodiscard]] int day() const noexcept { return this->day_; }

private:
    int year_ = 0;
    int month_ = 1;
    int day_ = 1;
};
\end{lstlisting}

这个类没有 setter。不是因为 setter 永远不好，而是因为单独改月份或日期可能瞬间破坏不变量。更合适的接口是业务动作，例如 \code{next_day()} 或 \code{add_days()}，由类内部维护规则。

\section{类背后的普通对象模型}

类不是魔法。可以把它理解成“数据字段加一组只能通过对象调用的函数”。下面的成员函数：

\begin{lstlisting}[language=C++,caption={成员函数}]
[[nodiscard]] bool contains(double value) const noexcept
{
    return this->low_ <= value && value <= this->high_;
}
\end{lstlisting}

近似等价于一个接收当前对象地址的普通函数：

\begin{lstlisting}[language=C++,caption={成员函数的近似展开}]
bool TemperatureRange_contains(const TemperatureRange* self,
                               double value) noexcept
{
    return self->low_ <= value && value <= self->high_;
}
\end{lstlisting}

真实 \code{private} 成员不能被外部普通函数访问，这里只是帮助理解：\code{this} 就是当前对象的入口。构造函数则是在对象创建时运行的一段初始化代码，用来把字段放到有效状态。析构函数是在对象离开生命周期时运行的清理代码。用这个普通代码模型看类，封装和生命周期会更具体。

\section{不变量让调用者少做防御}

如果 \code{Date} 保证自己永远有效，那么格式化函数就不需要每次重新检查月份范围：

\begin{lstlisting}[language=C++,caption={建立在不变量上的函数}]
std::string format_date(const Date& date)
{
    return std::format("{:04}-{:02}-{:02}",
                       date.year(),
                       date.month(),
                       date.day());
}
\end{lstlisting}

这个函数可以信任 \code{date.month()} 在 \code{1..12}，因为坏对象无法构造出来。不变量的价值就在这里：检查集中在边界，内部逻辑不必到处防御。

\section{测试 Date 的规则矩阵}

日期测试要覆盖规则，而不是只测 getter：

\begin{longtable}{p{0.25\linewidth}p{0.28\linewidth}p{0.35\linewidth}}
\toprule
场景 & 输入 & 期望 \\
\midrule
普通日期 & \code{2026, 6, 27} & 构造成功 \\
月份过小 & \code{2026, 0, 1} & 拒绝 \\
月份过大 & \code{2026, 13, 1} & 拒绝 \\
普通二月 & \code{2026, 2, 29} & 拒绝 \\
闰年二月 & \code{2024, 2, 29} & 构造成功 \\
日期为零 & \code{2026, 1, 0} & 拒绝 \\
\bottomrule
\end{longtable}

有了这个矩阵，类的职责就很清楚：任何公开可见的 \code{Date} 对象都必须是真实日历日期。
